\documentclass[11pt,a4paper]{article}

\usepackage[margin=2.2cm]{geometry}
\usepackage[T1]{fontenc}
\usepackage[utf8]{inputenc}
\usepackage{lmodern}
\usepackage{amsmath,amssymb}
\usepackage{enumitem}
\usepackage{parskip}

\begin{document}

\section*{Suggested Answers --- trial\_exam\_7}

\subsection*{Section 1 --- Multiple Choice (Q1--Q20)}
\begin{itemize}
    \item Q1 B
    \item Q2 A
    \item Q3 B
    \item Q4 C
    \item Q5 A
    \item Q6 A
    \item Q7 B
    \item Q8 B
    \item Q9 C
    \item Q10 A
    \item Q11 A
    \item Q12 C
    \item Q13 C
    \item Q14 A
    \item Q15 B
    \item Q16 B
    \item Q17 B
    \item Q18 B
    \item Q19 A
    \item Q20 A
\end{itemize}

\subsection*{Section 2 --- Short Questions (Q1--Q6)}
\begin{enumerate}[label=\arabic*.]
    \item \textbf{Symbol and non-instantaneous:} a symbol is the modulated waveform element used to carry information over a finite time interval; it is not instantaneous because it occupies a duration (not a mathematical point in time).
    \item \textbf{ISI key idea:} due to multipath/channel memory, delayed energy from previous symbols overlaps the current symbol at the receiver.
    \item \textbf{Spectral efficiency:} $\eta=\frac{R_b}{B}$.
    \item \textbf{Duplexing:} how uplink and downlink share resources (time division and/or frequency division).
    \item \textbf{Resource allocation key challenge:} interference management (interference depends on user activity/geometry and changes over time, making it hard to keep users separable while meeting QoS).
    \item \textbf{Increasing $M$:} bits per symbol increase ($\log_2(M)$), but typical SNR requirements also increase for the same error probability.
    
\end{enumerate}

\subsection*{Section 3 --- Long / Applied Questions}

\subsubsection*{Question 1 --- Link Budget and Feasibility}
Given: $f=950$ MHz, $P_t=20$ dBm, $G_t=1$ dBi, $G_r=2$ dBi, $d=3.2$ km, $N=-108$ dBm, $\gamma_{\min}=8$ dB.

\begin{enumerate}[label=\alph*.]
    \item \textbf{Free-space path loss:}
    \[
    L_{\mathrm{FSPL}}=32.44+20\log_{10}(950)+20\log_{10}(3.2)\approx 102.10\ \mathrm{dB}.
    \]
    \item \textbf{Received power:}
    \[
    P_r=P_t+G_t+G_r-L_{\mathrm{FSPL}}=20+1+2-102.10\approx -79.10\ \mathrm{dBm}.
    \]
    \item \textbf{Minimum required received power:}
    \[
    P_{\min}=N+\gamma_{\min}=-108+8=-100\ \mathrm{dBm}.
    \]
    \item \textbf{Feasible?} yes, since $P_r\approx -79.10 > -100$ dBm.
\end{enumerate}

\subsubsection*{Question 2 --- Indoor Path Loss}
Given: distance loss $60$ dB, concrete wall losses $10$ dB and $7$ dB, door $5$ dB, glass window $2$ dB, $P_t=8$ dBm, $G_t=G_r=0$ dBi, sensitivity $-75$ dBm.

\begin{enumerate}[label=\alph*.]
    \item Total path loss:
    \[
    L_{\text{total}}=60+10+7+5+2=84\ \mathrm{dB}.
    \]
    \item Received power:
    \[
    P_r=P_t-L_{\text{total}}=8-84=-76\ \mathrm{dBm}.
    \]
    \item Link works? $-76$ dBm is $1$ dB below $-75$ dBm sensitivity, so \textbf{no}.
    \item Two improvements (examples): increase transmit power; use higher-gain antennas/better placement; reduce obstacles (or add relay); use more sensitive receiver.
\end{enumerate}

\subsubsection*{Question 3 --- Frequency Reuse}
Given reuse factor $N=7$.

\begin{enumerate}[label=\alph*.]
    \item Spectrum fraction per cell:
    \[
    \frac{1}{7}\approx 0.143\ (14.3\%).
    \]
    \item For $N=10$:
    \[
    \frac{1}{10}=0.10\ (10\%).
    \]
    \item Impact of $N$:
    \begin{itemize}
        \item Larger $N$ reduces co-channel interference (cells are farther apart spectrally).
        \item Smaller $N$ increases bandwidth per cell and typically increases capacity per cell, but with higher co-channel interference.
    \end{itemize}
    
\end{enumerate}

\subsubsection*{Question 4 --- Multi-hop Routing}
Feasible links: $1\to 2,\ 1\to 3,\ 2\to 3,\ 2\to 4,\ 3\to 4,\ 1\to 4$.

Energies: $E_{12}=2.1,\ E_{13}=2.9,\ E_{23}=1.8,\ E_{24}=3.0,\ E_{34}=1.7,\ E_{14}=6.2$.

\begin{enumerate}[label=\alph*.]
    \item All feasible routes from node $1$ to node $4$:
    \begin{itemize}
        \item Route A: $1\to 4$
        \item Route B: $1\to 2\to 4$
        \item Route C: $1\to 3\to 4$
        \item Route D: $1\to 2\to 3\to 4$
    \end{itemize}
    \item Total energy per route:
    \[
    E_A=E_{14}=6.2
    \]
    \[
    E_B=E_{12}+E_{24}=2.1+3.0=5.1
    \]
    \[
    E_C=E_{13}+E_{34}=2.9+1.7=4.6
    \]
    \[
    E_D=E_{12}+E_{23}+E_{34}=2.1+1.8+1.7=5.6
    \]
    \item Minimum-energy route: Route C ($1\to 3\to 4$) with $E_{\min}=4.6$.
    \item Advantage: better feasibility/coverage over long distances. Disadvantage: increased delay/overhead due to multiple hops.
\end{enumerate}

\subsubsection*{Question 5 --- Modulation and Throughput}
Bandwidth: $B=140$ kHz $=1.40\times 10^5$ Hz.

\begin{enumerate}[label=\alph*.]
    \item Shannon capacity at SNR $=10$ dB:
    \[
    C=B\log_2(1+\mathrm{SNR})\approx 4.84\times 10^5\ \text{bit/s}\ (\approx 484.32\ \text{kbps}).
    \]
    \item Highest usable scheme at $10$ dB: scheme-4 (threshold $\ge 8$ dB). Throughput $\approx 2.0$ bit/s/Hz.
    \[
    R\approx 2.0\cdot B=2.0\cdot 140000=2.80\times 10^5\ \text{bit/s}\ (280\ \text{kbps}).
    \]
    \item From scheme-5 (at 14 dB) to scheme-4:
    scheme-5 is usable for SNR $\ge 13$ dB, so
    \[
    \Delta \approx 14-13=1\ \text{dB}.
    \]
\end{enumerate}

\subsubsection*{Question 6 --- Coverage Planning}
\begin{enumerate}[label=\alph*.]
    \item \textbf{Single-site:} place the base station at/near the hill crest or highest elevation for improved diffraction/coverage toward both towns.
    \item \textbf{Two-site (base + relay):} place the relay on the far-side of the hill so the shadowed town gets a stronger hop.
    \item \textbf{Compare:} two-site improves coverage and reliability in shadowed regions but costs more (extra site/backhaul and coordination).
    \item \textbf{Most important effects:} hill shadowing/blockage, diffraction around the obstacle, and free-space path loss with distance (plus multipath/reflections).
    
\end{enumerate}

\end{document}

